% Generated by scripts/build_paper_tables.py; do not edit by hand.
\par
\begin{table}[htbp]
\centering
\caption{确定性价格感知反馈与原反馈的费用比较（2—12月，元）}\label{tab:mpc-comparison}
\footnotesize
\setlength{\tabcolsep}{4pt}
\renewcommand{\arraystretch}{1.15}
\begin{tabular}{lrrrr}
\toprule
方案 & 反馈规则 & 常规购电费 & 紧急购电费 & 总费用 \\
\midrule
问题2 & 贪心 & 13,291,625.72 & 616,550.23 & 13,908,175.95 \\
问题2 & 确定性滚动 & 13,287,773.99 & 654,732.22 & 13,942,506.21 \\
问题3 & 贪心 & 13,177,405.10 & 138,430.89 & 13,315,835.99 \\
问题3 & 确定性滚动 & 13,062,906.94 & 466,563.01 & 13,529,469.95 \\
问题4-2 & 贪心 & 14,032,328.84 & 643,516.03 & 14,675,844.87 \\
问题4-2 & 确定性滚动 & 14,028,077.93 & 681,762.73 & 14,709,840.66 \\
问题4-3 & 贪心 & 13,917,214.62 & 127,670.62 & 14,044,885.25 \\
问题4-3 & 确定性滚动 & 13,807,269.57 & 451,105.26 & 14,258,374.82 \\
\bottomrule
\end{tabular}
\par\smallskip
\begin{minipage}{0.97\linewidth}\footnotesize 常规购电费为原计划费与调整费之和。两种反馈共用参数、发布时间、1月热启动和正式期初库存；反馈改变库存后，后续常规计划也随之改变。\end{minipage}
\end{table}
\par
